% Research notes for the 2026-10-12 continuation.  This is a section fragment.
% Citation keys expected: AHT2006, Lackenby2021, Lackenby2026.
\section{A finite interface for exponentially subdivided bundle boundaries}
\label{sec:port-boundary-theorem}

The all-dimensional parallelity construction already provides the base
surface of each source bundle as families of polygons, together with the
set of core-facing sides, called ports.  A base profile records Euler
characteristic, orientability, boundary count, and the set of ports in its
component.  It does not record which boundary circle contains a port or
the cyclic order of the ports.  The polygon packet contains that information,
but reading it by walking a boundary circle may require a number of steps
proportional to the normal weight.  We give a reduction that extracts the
needed finite interface directly from the compressed packet.

This is an implementation-oriented specialization of earlier algorithmic
topology.  Lackenby's Theorems 9.2 and 9.3 in \cite{Lackenby2021} already
give polynomial algorithms for parallelity bundles with attachment data;
the explanation of the output explicitly includes cyclic order around
vertical annuli and coherent horizontal-sheet labels for product bundles.
We do not claim a new existence theorem for that general task.  The
contribution here is an explicit endpoint-matching reduction for the
repository's all-dimensional polygon packet, its source-labelled fibre
charts, and the safe distinction between full and partial disc attachments.
The repository also already contains the degree-two ordered-boundary
principle in Report 68, summarized in
\texttt{synthesis/completion\_theorems.tex}: endpoint count, sum, square
sum and length recover the two endpoint occurrences of each cut path.
The present construction integrates that principle with the later complete
parallelity packet and adds zero-gap flags, the signed fibre cover, and
executable source-interface verification.

\subsection{The abstract input and the exact output}

Let $B$ be a finite compact surface assembled from a finite, possibly very
large, collection of polygon copies.  Assume that copies, side pairings and
their endpoint maps are represented by $s$ binary interval records on a
corner universe $\{0,\ldots,N-1\}$.  Pairings are integral affine
isometries, possibly order reversing.  Each polygon-side identification
gives identifications between its endpoint corners.  Let $V$ denote the
resulting corner-pairing system.  Its orbits are precisely the vertices of
the quotient polygonal surface.  The input is required to describe a
surface; this is a geometric hypothesis, rather than a conclusion of the
following orbit computation.

Distinguish $p$ individual unglued sides, each specified by one binary
rank and a unique port label.  Write $e_{j,0}$ and $e_{j,1}$ for the two
\emph{endpoint flags} of port $j$.  Endpoint flags remain distinct even
when their corner addresses represent the same quotient vertex.
Every other unglued side is an old-boundary side.  Let $R$ be $V$ together
with the endpoint pairing of every old-boundary side.  Crucially, the
endpoint pairings of the ports are omitted from $R$.

The desired output consists of the following data:
\begin{enumerate}
\item a matching $\tau$ of the $2p$ endpoint flags, pairing the two ends of
  each component of $\partial B$ after the open port sides are deleted;
\item a bit $z$ for every matched pair, specifying whether the intervening
  component is a single quotient vertex, rather than an arc containing
  old-boundary sides;
\item the cyclic lists of labelled and directed ports on every boundary
  circle that contains a port, with the $z$ bits on the intervening gaps.
\end{enumerate}
If a fibre-reversal bit is attached to each polygon gluing, we also recover
the fibre transport along each intervening component.  The cyclic lists
are canonically chosen up to reversal; this does not choose a common
orientation for different boundary components of the same base surface.

\begin{theorem}[Marked boundary extraction]
\label{thm:marked-boundary}
The matching $\tau$, the gap bits $z$, and the cyclic port lists can be
computed with $O(p^2)$ calls to ordinary interval-orbit counting on systems
with $O(s+p)$ records and universe size $N$.  Consequently they can be
computed in time polynomial in $s+p+\log(N+1)$ by the Agol--Hass--Thurston
orbit-counting theorem.  Their combinatorial output size is $O(p)$.

For the source parallelity packet of a triangulation with $T$ tetrahedra
and coordinate bit length $L$, this is polynomial in $T+L$, since
$s=O(T)$, $p\leq66T$, and $\log(N+1)=O(L+\log(T+1))$.
\end{theorem}
\begin{proof}
First quotient the corner universe by $V$.  Interior vertices are
irrelevant to $\partial B$.  Boundary vertices and the unglued sides form
a finite cellulation of the disjoint union of circles $\partial B$.
Deleting the open interiors of the port sides leaves circles with no
marked endpoints and closed intervals with exactly two marked endpoint
flags.  An interval is allowed to have no side and consist of a single
vertex.  This happens precisely when two consecutive port sides share
a quotient vertex.  If one port itself has coincident endpoint vertices,
its two endpoint flags must still be counted separately.

Adding the old-boundary side pairings to $V$ identifies exactly the
vertices within these remaining components.  Hence every $R$-orbit
containing a marked flag contains exactly two marked flags.  Its two flags
give $\tau$.  Notice that the assertion concerns flags, not distinct
integer addresses: two marked flags may have the same address.

Here is a reduction to unweighted orbit counting that avoids any assumption
about the performance of a weighted implementation.  For a pairing system
$Q$, write $c(Q)$ for its number of orbits.  For two corner addresses $x,y$,
adjoin the singleton pairing $x\leftrightarrow y$.  Then
\[
 c\bigl(Q\cup\{x\leftrightarrow y\}\bigr)
 =\begin{cases}
 c(Q),&x\sim_Q y,\\
 c(Q)-1,&x\not\sim_Q y.
 \end{cases}
\]
The reason is that one added edge merges exactly two different connected
components, or merges none.  Compute $c(R)$ once and apply this test to
each pair of marked addresses.  The resulting equivalence classes on
endpoint flags have size two, giving $\tau$ in $O(p^2)$ counting calls.
Repeated addresses cause no difficulty.

For each pair matched by $\tau$, apply the same test to $V$.  The two flags
represent the same quotient vertex if and only if the remaining interval
has no old-boundary side.  Indeed, an interval with distinct ends cannot
have its ends identified in a cellulation of a one-manifold.  This gives
the bits $z$ with at most $p$ additional counting calls.

Let $\sigma(e_{j,0})=e_{j,1}$ and
$\sigma(e_{j,1})=e_{j,0}$.  The multigraph on endpoint flags whose edges
are the $\sigma$ pairs and $\tau$ pairs is 2-regular.  Its components are
exactly the boundary circles meeting ports.  Traverse alternating port
and gap edges to obtain each cyclic list.  Each flag and port is visited
only a constant number of times.  Canonicalize by beginning with the
smallest label and choosing the lexicographically smaller direction.
Thus the postprocessing and output size are $O(p)$, apart from the usual
sorting of labels.  The claimed polynomial bound follows from the basic
AHT theorem \cite{AHT2006}.
\end{proof}

This theorem gives a polynomial algorithm as an abstract reduction.
Using the repository's current orbit engine supplies an implementation
whose certificate replay is bounded by its realized traces.  Identifying
that specific implementation with the polynomial AHT theorem requires
its own complexity proof; the theorem above does not silently supply it.

\subsection{One-query weighted realization}

There is a faster practical realization using the existing weighted
histogram primitive.  Give endpoint flag $e_{j,a}$ the singleton weight
$2^{2j+a}$ at its corner address and give all other points zero weight.
Overlapping singleton weights are added.  Every positive-weight $R$-orbit
then has a weight containing exactly two set bits.  Those two bits recover
the matched flags.  The uniqueness of the endpoint labels prevents carries
and implies that a positive mask has multiplicity one.  Thus the positive
part of the histogram has exactly $p$ rows; arbitrarily many unmarked
components contribute only to the zero-weight row.  A second query on $V$
gives the endpoint vertex partition and the zero-gap information.

The scalar weights have $2p$ bits, not $2^p$ bits.  In particular, the
binary encoding remains polynomial in $T+L$ even when a boundary path
contains exponentially many old-boundary sides.  This is a bounded-mark
use of the weighted primitive; it does not require a general bound on
all possible weighted histograms.

\subsection{Transport of the interval fibre}

Suppose the polygon packet is equipped with compatible local interval
fibres.  For every gluing, let $\epsilon\in\{0,1\}$ record reversal of
the fibre.  Lift its corner maps to two sheets by
\[
                (x,a)\longmapsto(g(x),a\mathbin\oplus\epsilon).
\]
An old-boundary edge has zero fibre reversal in its polygon chart.  Let
$\widetilde R$ be the resulting lifted system on $2N$ corner addresses.
Each marked endpoint has two lifted flags.

\begin{proposition}[Signed gap extraction]
\label{prop:signed-gap}
The lifted endpoint matching determines one bit $\eta(e,f)$ for every
gap matched by $\tau$: continuation from sheet zero at $e$ reaches sheet
$\eta(e,f)$ at $f$.  It is computable in polynomial encoded time by the
same orbit-count reduction, or by one weighted query with $4p$ endpoint
bits.  A valid interval bundle has precisely two lifts of each gap, so the
lifted matching commutes with the deck involution.

If the total interval bundle is orientable, the xor of the gap bits around
every cyclic port list is zero.
\end{proposition}
\begin{proof}
Every gap is an interval or a point, and an interval bundle over either
is trivial.  Its two sheets therefore remain distinct.  The lifted
components pair the two lifts of each endpoint bijectively with the two
lifts of the other endpoint.  This bijection is either the identity or
the transposition, and defines $\eta$.  The same orbit membership test
as in Theorem~\ref{thm:marked-boundary} computes it.

For the last assertion, the orientation character of the vertical
line bundle equals that of $B$, because their sum is the orientation
character of the total space.  The restriction of $TB$ to a boundary
circle splits into the tangent line of the circle and its outward normal
line, both orientable.  Hence the fibre character is trivial on every
boundary circle.  Traversal through a port has zero fibre change in its
local polygon chart, so the xor of the recorded gap changes is exactly
this zero character.  The same reasoning applies to a boundary circle
of a nonorientable base whose total interval bundle is orientable.
\end{proof}

The zero-holonomy condition is a useful local consistency check.  It does
not imply that the full bundle is a product: an orientable interval bundle
over a M\"obius band has trivial fibre holonomy on its boundary but has
nontrivial fibre holonomy on its core curve.

\subsection{What is and is not determined}

The new interface determines the partition of ports into boundary circles,
the direction of each port relative to its endpoint chart, their cyclic
order, whether consecutive ports actually touch, and the fibre transport
between them.  It does not enumerate the old-boundary sides.  If a base
profile has $b$ boundary components and $c$ recovered marked circles, the
remaining $b-c$ unmarked circles are retained as a binary count.

The interface is a complete description of the \emph{marked vertical
boundary}, in its supplied endpoint and fibre charts.  It is not by itself
a complete classification of the interval bundle relative to all markings
on its horizontal boundary.  Such markings, as well as coherent orientation
choices between different boundary circles, must be transported separately
when required by a patterned hierarchy.

The omission of cyclic order from a base profile is a real loss of
information.  A disc with four labelled boundary arcs in order
$(a,b,c,d)$ is not homeomorphic, relative to the labels, to a disc with
order $(a,c,b,d)$.  A boundary homeomorphism preserves cyclic order up to
reversal.  Both discs have identical Euler characteristic, boundary
count, orientability and port set.  Thus no algorithm using only the
current numerical profile can construct the required marked boundary
homeomorphism in general.

\section{Safe replacement of disc-based source bundles}
\label{sec:disc-bundle-replacement}

Let $K$ be a compact 3-manifold with corners and let
$E=D^2\times I$ be a product component of the source parallelity bundle.
Assume that
\[
              E\cap K=A\times I,
              \qquad A\subseteq\partial D^2,
\]
where $A$ is a union of the port sides, with the given source attachment
charts.  The interiors of these interfaces are embedded and different
bundle components have disjoint interfaces.  Consecutive port sides may
meet; the cyclic interface tells us when to combine them.  These are
topological gluing hypotheses already supplied by the full source packet.

\begin{theorem}[Disc bundle replacement with partial attachment]
\label{thm:disc-replacement}
There are exactly three possibilities.
\begin{enumerate}
\item If $A=\varnothing$, then $E$ is a separate 3-ball component.
\item If $A=\partial D^2$, then adjoining $E$ is precisely the attachment
  of one 2-handle, with attaching annulus $\partial D^2\times I$.
\item If $A$ is nonempty and proper, let $r$ be its number of connected
  components.  Each component is a closed interval.  Up to transport
  through collars of the attaching discs, adjoining $E$ is the addition
  of $r-1$ relative 1-handles to $K$.
\end{enumerate}
The homeomorphism preserves any explicitly transported boundary markings.
In the third case it can be chosen fixed outside collars of the attaching
discs.  It need not fix the entire originally embedded $K$ pointwise.

For all attached disc-base components together, the number of relative
handles added in cases (2) and (3) is at most the total number of ports.
Unattached disc-base components may be retained by binary multiplicity.
\end{theorem}
\begin{proof}
The first two cases are immediate from the product description.  In case
(3), write $A=A_1\cup\cdots\cup A_r$ in cyclic order.  There is a set of
$r-1$ disjoint properly embedded arcs in the disc, with endpoints outside
$A$, which cuts the disc into $r$ discs, each containing exactly one
$A_i$ on its boundary.  One construction successively cuts off $A_1$,
then $A_2$, and so on, choosing endpoints in the two adjacent complementary
boundary intervals.  The dual graph of these cuts is a tree.

Take the product of these arcs with $I$.  They are $r-1$ disjoint vertical
discs splitting $E$ into $r$ balls.  Each ball meets $K$ in just the disc
$A_i\times I$.  Absorb it into a collar of that disc in $K$.  This gives
a 3-manifold $K'$ homeomorphic to $K$ by a homeomorphism fixed away from
the attaching collars.  Gluing the $r$ pieces back across the cut discs
is equivalently accomplished by attaching $r-1$ 1-handles between their
collar patches.  Product charts on the cut discs give the attaching maps.
Boundary markings on the pieces are carried by the same homeomorphisms,
so they are not discarded by the replacement.

If a disc component uses $p_i$ ports, then a full-circle attachment uses
one new 2-handle and has $p_i\geq1$.  A partial attachment has
$r\leq p_i$ and uses $r-1\leq p_i$ new 1-handles.  Summing gives the bound.
The components with no ports are disjoint balls and may occur with large
multiplicity, so a compressed representation must count rather than
explicitly list them.
\end{proof}

The case $r=1$ is particularly useful: the whole disc bundle is a collar
addition and introduces no relative handle.  Replacing every disc base
indiscriminately by a 2-handle would be wrong for such partial attachments.
For example, with $r=2$ the correct operation adds a 1-handle, whereas a
2-handle would have the opposite Euler-characteristic contribution.

\begin{corollary}[Computable relative Euler contribution]
For a disc product component,
\[
 \chi(K\cup E)-\chi(K)=
 \begin{cases}
 1,&A=\varnothing\text{ or }A=\partial D^2,\\
 1-r,&A\neq\varnothing,\ A\neq\partial D^2.
 \end{cases}
\]
The cyclic interface determines the applicable case and $r$ without
expanding the polygon packet.  On a marked boundary circle, $A$ is the
entire circle exactly when every gap has $z=1$.  Otherwise $r$ is the
number of gaps with $z=0$.
\end{corollary}
\begin{proof}
Euler characteristic under gluing gives
$\chi(K\cup E)-\chi(K)=\chi(E)-\chi(A\times I)=1-\chi(A)$.
The asserted values and the description by gaps follow immediately.
\end{proof}

This is a local simplification theorem and a bounded relative handle
presentation.  It does not prove that the transported boundary pattern is
essential, that the resulting handle structure is admissible, or that its
0-handles remain in a fixed finite list of types through a hierarchy.
Those are distinct hypotheses in the hierarchy theorems
\cite{Lackenby2026}.  In particular, reordering and refining the relative
handle additions into the hierarchy's required source handle structure
needs an explicit construction and its own size bound.  The theorem also
does not authorize discarding unmarked ball components when their inherited
horizontal boundary markings still have to be checked.

\subsection{Certified implementation and finite geometric comparison}

The supplemental prototype performs three marked weighted-orbit queries:
the vertex partition, the gap matching, and its fibre double cover.  Its
independent checker reconstructs the query systems from the authenticated
source polygon packet and replays all three weighted certificates.  It
does not invoke the producer's query builder, mask decoder, or cyclic
traversal.  Source authentication itself remains the responsibility of
the maintained \texttt{verify\_normal\_cut\_parallelity} operation.

The reproducible focused audit records 11 successful test methods,
184 finite geometric comparisons, and 54 authenticated source comparisons.
The finite oracle first glues actual polygon corners, checks their fibre
charts, and walks the actual boundary cellulation; it does not use the
compressed interval-query construction.  Among the synthetic cases are
all 126 subsets of marked sides of discs with one through six sides,
including the monogon case where the two flags of one port have the same
corner address.  Annular and M\"obius base charts test fibre transport.
Source cases include relabellings, interior handles, and the existing rich
3-manifold fixtures.

Two source examples expose the distinction needed by
Theorem~\ref{thm:disc-replacement}.  In the two-tetrahedron source
\texttt{fibonacci\_lst\_02}, coordinate rows
$(0,0,1,1,1,0,0)$ and $(1,1,1,1,0,0,0)$ give a single disc base supported
by one selected 2-handle and no selected 0- or 1-handles.  Its eight ports
form one cyclic list with no old-boundary gaps, so it is precisely a
full-circle 2-handle attachment.  On the other hand, three parallel
meridians in the two-tetrahedron layered torus give two attached disc
bases, each with one proper core interval.  They are collar additions,
introducing no relative handle, together with two unattached ball products.
Complete source and interface certificates for these examples accompany
the prototype.

The large-coordinate regression uses $2^{8192}$ parallel meridians in a
four-tetrahedron source.  It checks exact JSON transport and independent
replay with producer routines disabled.  The literal oracle refuses this
input because it exceeds its explicit expansion cap.  Hostile tests alter
gap adjacency, cyclic directions, fibre bits, source digests and weight
masks.  Shared discovery limits, shared verification limits, and cancellation
are checked separately.  Exhaustion returns no geometric interface.

These observations validate the additional interface computation on the
tested inputs.  They do not constitute a proof of a global hierarchy
bound, or a replacement for the topological hypotheses in the preceding
theorems.  In particular, the returned disc-replacement record is a local
relative-handle plan; the prototype performs no surgery.

\subsection{Research priorities exposed by this reduction}

The finite interface suggests several concrete next questions.
\begin{enumerate}
\item Construct a certificate for the marked boundary homeomorphism and
  the induced corner pattern after each disc-bundle replacement.  The
  certificate should use the cyclic interface and source charts, rather
  than the expanded base disc.
\item Extend marked-boundary extraction to the endpoints and vertices of
  a compressed boundary pattern.  Determine conditions under which a
  polynomial number of pattern events is preserved by a hierarchy step.
\item Produce a polynomial conversion from the resulting relative
  handle presentation to an admissible handle structure of uniform type.
  The inherited $30T$ core bound alone does not establish this.
\item For annular and nonorientable bases, retain the full relative
  orientation character between marked boundary circles.  Determine the
  least extra finite data needed for a bundle homeomorphism relative to
  both its core interfaces and its horizontal pattern.
\item Prove the necessary polynomial bound for the repository's actual
  interval-orbit implementation, or introduce a certified implementation
  that follows an already proved AHT reduction bound.
\end{enumerate}
